\documentclass[UTF8,11pt]{ctexart}
\usepackage[a4paper,margin=2.5cm]{geometry}
\usepackage{amsmath,amssymb,amsthm,mathtools}
\usepackage{booktabs,longtable,hyperref}

\newtheorem{definition}{定义}[section]
\newtheorem{lemma}[definition]{引理}
\newtheorem{theorem}[definition]{定理}
\newtheorem{corollary}[definition]{推论}
\newtheorem{proposition}[definition]{命题}
\newtheorem{remark}[definition]{说明}

\title{IC2 2.8.221 Mark I 计算优化的形式化证明}
\author{IC2 Nuclear Reactor Optimizer}
\date{\today}

\begin{document}
\maketitle
\tableofcontents

\section{范围与假设}

本文只讨论 IC2 Experimental 2.8.221 的 EU 反应堆、铀燃料和项目当前支持的组件集。反应堆使用两阶段、行优先、立即写回结算。Mark I 优化器开启自动续棒，排行榜目标严格为平均发电功率 EU/t。

\section{状态转移与固定点}

\begin{definition}[约化热状态]
令
\[
S_n=(H_n,L_n,C_n,D_n),
\]
其中 $H_n$ 为堆体热量，$L_n$ 为槽位组件 ID，$C_n$ 为组件热量，$D_n$ 为非燃料组件耐久。自动续棒时燃料耐久不进入约化状态。
\end{definition}

固定布局定义确定性状态转移
\[
S_{n+1}=F(S_n).
\]

\begin{proposition}[燃料产热与温度无关]
一根内部燃料棒具有 $p$ 个有效脉冲时，产热为
\[
Q(p)=4\sum_{i=1}^{p}i=2p(p+1).
\]
有效脉冲数只由内部脉冲、相邻燃料和相邻有效反射板决定，因此温度不直接参与产热。
\end{proposition}

\begin{theorem}[固定点快进]
若未发生临界、组件移除或融毁，且 $F(S_n)=S_n$，则对任意 $m\ge1$，有
\[
F^m(S_n)=S_n.
\]
\end{theorem}

\begin{proof}
对 $m$ 使用数学归纳法。$m=1$ 由条件成立。若 $F^k(S_n)=S_n$，则
\[
F^{k+1}(S_n)=F(F^k(S_n))=F(S_n)=S_n.
\]
燃料耐久不改变约化状态转移，自动续棒保持燃料类型和脉冲拓扑，因此每周期 EU/t 也相同。
\end{proof}

\section{发电骨架与冷却赋值}

\begin{definition}[组件划分]
令
\[
\mathcal P=\{\text{燃料、反射板}\},\qquad
\mathcal C=\{\text{散热、换热、储热、冷凝、隔板}\}.
\]
两者互不相交。
\end{definition}

\begin{definition}[二层表示]
对完整布局 $L$，$P(L)$ 保留 $\mathcal P$ 中的组件并把其他槽位变为空槽；$C(L)$ 记录 $P(L)$ 空位中的冷却组件赋值。
\end{definition}

\begin{theorem}[二层表示双射]
映射
\[
\Phi:L\longmapsto(P(L),C(L))
\]
是完整库存有效布局集合与合法二层表示集合之间的双射。
\end{theorem}

\begin{proof}
若 $L_1\ne L_2$，至少有一个槽位不同。该差异若属于 $\mathcal P$，则两个骨架不同；否则两个冷却赋值不同，故 $\Phi$ 为单射。反之，给定合法 $P$ 和 $C$，把 $C$ 唯一写回 $P$ 的空位即可重建一个完整布局，故 $\Phi$ 为满射。
\end{proof}

\begin{corollary}[库存约束保持]
燃料、燃料棒总数和反射板约束只作用于 $P$，冷却组件上限只作用于 $C$。由于两类组件不相交，完整布局满足库存约束当且仅当两层分别满足对应约束。
\end{corollary}

\section{骨架功率与安全剪枝}

\begin{theorem}[Mark I 骨架功率精确性]
设第 $j$ 个燃料组件有 $r_j$ 根内部棒、$a_j$ 个内部脉冲和 $d_j$ 个相邻燃料或反射板，则在组件不移除的 Mark I-I 布局中
\[
P=5\sum_j r_j(a_j+d_j)\quad\mathrm{EU/t}.
\]
具有相同发电骨架的所有 Mark I-I 冷却完成布局拥有相同 EU/t。
\end{theorem}

\begin{proof}
能量阶段的每个有效脉冲产生 5 EU/t。冷却组件不接收或返回铀脉冲。Mark I-I 中组件不被移除，因此脉冲拓扑保持不变。
\end{proof}

\begin{lemma}[活动有限反射板排除]
与燃料相邻的有限反射板每周期至少增加一次耐久损耗，最终必然移除。因此它不可能无限期满足 Mark I-I 的无组件损坏条件。
\end{lemma}

\begin{theorem}[可配置 Top-$K$ 排行榜下界剪枝]
令 $K=\texttt{result\_limit}$。若当前存在 $K$ 个不同规范布局组，其功率均不小于 $\tau_K$，则任何骨架功率上界 $U<\tau_K$ 的子树都不可能包含 Top-$K$ 候选。若 $K=1$ 且只要求任意一个全局最优见证，则条件可加强为 $U\le\tau_1$。
\end{theorem}

\begin{proof}
由骨架功率精确性和部分骨架上界的乐观性，该子树中任一 Mark I-I 布局的功率 $p\le U$。已有 $K$ 个不同规范组的功率至少为 $\tau_K>U$，故该布局不可能进入 Top-$K$。当 $K>1$ 时算法不剪枝 $U=\tau_K$，从而保留并列候选；当 $K=1$ 时，首个功率为 $\tau_1$ 的见证已经保留，$U\le\tau_1$ 的后代不可能严格改善目标。
\end{proof}

\begin{theorem}[不会遗漏优解]
二层生成器不会遗漏任何可能进入所请求 Mark I Top-$K$ 的库存有效布局；Top-1 至少保留一个全局最大功率见证。
\end{theorem}

\begin{proof}
由双射定理，每个完整布局属于唯一骨架和冷却赋值。活动有限反射板骨架由引理证明不可行；不能改善第 $K$ 名门槛的骨架由排行榜下界定理安全关闭；其余骨架的冷却赋值全部生成并执行官方热学验证。因此所有可能进入 Top-$K$ 的布局均被验证。Top-1 的等号剪枝发生前已经存在同功率见证，故不会失去最优目标值。
\end{proof}

所谓“先找第一个”应在一次遍历内实现：$K=1$ 时首个候选立即建立共享门槛，后续分支用上界关闭。若为第一、第二、第三名各自从根重跑一次，则公共前缀被重复访问，正确但通常不具有效率优势。

\section{冷却子树计数}

设骨架有 $r$ 个空位，冷却组件上限为 $c_1,\ldots,c_k$。选择数量 $n_i$ 后，有标签赋值数为
\[
\frac{r!}{(r-\sum_i n_i)!\prod_i n_i!}.
\]
总数为
\[
N_C(r;c_1,\ldots,c_k)=
\sum_{\substack{0\le n_i\le c_i\\\sum_i n_i\le r}}
\frac{r!}{(r-\sum_i n_i)!\prod_i n_i!}.
\]
被证明排除的骨架直接把 $N_C$ 加入检查数和数学跳过数，无需构造各个完整布局。

\section{性能边界}

令 $N$ 为完整布局数，$K$ 为基础生成器实际构造的完整布局数，$M$ 为最终验证数，$B$ 为批量跳过数，则
\[
N=K+B,\qquad M\le K\le N.
\]
正常完成时
\[
\texttt{checked}=\texttt{evaluated}+\texttt{pruned}=N.
\]

设骨架成本、完整布局构造成本、最终验证成本、见证短证明成本和额外调度成本分别为 $C_s,C_c,C_t,C_w,C_o$，构造器尝试 $W$ 个见证。近似有
\[
T_{\mathrm{old}}=N(C_c+C_t),
\]
\[
T_{\mathrm{new}}=SC_s+KC_c+MC_t+WC_w+C_o.
\]
严格加速条件为
\[
(N-K)C_c+(N-M)C_t>SC_s+WC_w+C_o.
\]
所以可以证明最终验证不重复并保持完整计数，但见证失败会增加 $WC_w$，不能无条件证明所有输入的墙钟时间严格缩短。

\section{证书缓存}

缓存键由完整有标签布局、列数和最大模拟周期组成。键完全相同意味着确定性状态转移和停止条件相同，可以复用结果。镜像布局不共享证书，因为其行优先状态转移可能不同。缓存采用容量上限，避免无限增长。

\section{简单固定点证书与冷却见证构造器}

当前实现的 \texttt{prove\_simple\_fixed\_point} 从零热状态出发，对不含换热器、冷却单元和冷凝模块的简单布局最多执行 4 次生产状态转移。若安全可达状态 $x$ 满足 $F_L(x)=x$，则确定性给出对任意 $n\ge1$ 都有 $F_L^n(x)=x$，故可直接返回 Mark I-I 见证。证书失败只返回 \texttt{None}，统一评估入口随后调用原精确模拟器。

\texttt{construct\_simple\_cooling\_candidates} 按精确燃料产热和邻接关系放置库存内的直接散热组件，并尝试不同的行优先方向。构造结果必须通过上述可达固定点证书；未构造或未证明的情况一律回退完整冷却生成。

\begin{theorem}[构造快速通道不降低完整性]
在构造失败和证书未知时保留原完整生成；构造成功只增加一个已证明可行见证。于是 Top-$K$ 的可行候选集合不会减少。Top-1 中，见证实现完整骨架的精确功率后，跳过该骨架其余冷却完成不会损失更高功率解。
\end{theorem}

\begin{proof}
任何未由构造器找到的可行冷却完成仍由原生成器访问。证书成功的布局是真实可行布局，不是假阳性。对同一完整骨架 $s$ 的任意冷却完成 $c$，有 $P(s\oplus c)=P(s)$；Top-1 已保留一个取得 $P(s)$ 的见证后，其余完成至多并列。Top-$K$ 在 $K>1$ 时不应用这一提前停止，故并列规范组仍被生成。
\end{proof}

\section{构造式精确求解器的程序契约}

上一节的简单直连冷却见证构造器和本节的前沿动态规划、精确功率优先队列均已实现；通用外部约束合成器仍是下一阶段设计。热学证书无法判定时，当前二层冷却生成器及精确模拟回退仍是完整性基线。

给定验证边界 $T$，令 $\operatorname{Accept}_T(L)$ 表示当前精确模拟器从零热状态出发，在边界 $T$ 内把布局 $L$ 判定为 Mark I。任何松弛模型的 \texttt{SAT} 结果只能作为候选；\texttt{UNKNOWN}、超时或取消必须保留为未决，并最终回退到精确完整验证。

\section{骨架功率、产热与可实现性}

\begin{theorem}[骨架同时决定功率与产热]
设第 $j$ 个燃料组件具有 $r_j$ 根内部棒、$a_j$ 个内部脉冲、$d_j$ 个相邻燃料或有效反射板，并令 $q_j=a_j+d_j$。只要骨架保持存在，则
\[
P(s)=5\sum_j r_jq_j,
\qquad
Q(s)=2\sum_j r_jq_j(q_j+1).
\]
冷却完成不改变 $P(s)$ 或 $Q(s)$。
\end{theorem}

\begin{proof}
每根内部棒贡献 $q_j$ 个能量脉冲，每个脉冲为 5 EU/t，得到第一式。每根内部棒的产热为 $4\sum_{i=1}^{q_j}i=2q_j(q_j+1)$，乘以 $r_j$ 并求和得到第二式。冷却组件不接收铀脉冲，故不出现在两式中。
\end{proof}

令 $\mathcal S$ 为库存有效骨架集合，$\mathcal C(s)$ 为骨架 $s$ 的合法冷却完成集合，并定义
\[
\operatorname{Feas}_T(s)
\iff
\exists c\in\mathcal C(s),\quad
\operatorname{Accept}_T(s\oplus c).
\]
于是
\[
\operatorname{OPT}_T=
\max_{s\in\mathcal S:\operatorname{Feas}_T(s)}P(s),
\qquad
P_{\max}=\max_{s\in\mathcal S}P(s).
\]

\begin{theorem}[原始最大骨架定理]
$\operatorname{OPT}_T=P_{\max}$ 当且仅当至少存在一个功率为 $P_{\max}$ 的骨架满足 $\operatorname{Feas}_T$。
\end{theorem}

\begin{proof}
若存在这样的骨架及冷却完成，它达到全部骨架的统一上界，故全局最优。反之，若最优值等于 $P_{\max}$，按最大值定义必有一个可实现骨架取得该值。
\end{proof}

该条件不能省略。例如单铀、铱反射板和普通散热片各一个时，最大骨架具有 10 EU/t 和每周期 12 热，而唯一散热片持续外排不超过 6；有限总储热必然增长。去掉反射板后的骨架只有 5 EU/t 和 4 热，却可以形成 Mark I-I。

\section{可实现性与排除证书}

\begin{theorem}[可达周期证书]
对完整布局 $L$，若从初始状态沿安全前缀可达 $S$，存在 $k\ge1$ 使
\[
F_L^k(S)=S,
\]
且这 $k$ 个周期内没有临界、非燃料组件损坏或融毁，则该安全轨迹可以无限重复。
\end{theorem}

\begin{proof}
由确定性，对任意 $m\ge0$ 有 $F_L^{mk}(S)=S$，每个中间状态也逐段重复。因此已经验证的周期安全谓词对所有未来成立。
\end{proof}

代数固定点必须由精确模拟证明从零热状态可达。单周期固定点只是充分条件；若只检查它而不保留精确回退，会遗漏先升温后稳定或周期大于 1 的布局。

\begin{theorem}[持续热守恒排除]
在没有组件移除的安全轨迹中，令 $E_n$ 为总储热，$W_n$ 为实际外排热量，则
\[
E_{n+1}=E_n+Q(s)-W_n.
\]
若对所有周期有 $W_n\le V$ 且 $Q(s)>V$，则该布局不可能是 Mark I。
\end{theorem}

\begin{proof}
$E_n\ge E_0+n(Q(s)-V)$ 线性无界增长；但堆体和组件在不临界、不损坏条件下的总储热有有限上界，矛盾。
\end{proof}

对未完成骨架计算乐观持续外排上界 $\overline V(s)$：自散热片最多贡献自身 \texttt{self\_vent}，元件散热片最多贡献 \texttt{side\_vent} 乘以实际相邻可储热组件数；抽堆热和换热只搬运热量，储热、冷凝和隔板只增加有限容量。因此
\[
Q(s)>\overline V(s)
\Longrightarrow
\neg\operatorname{Feas}_T(s).
\]
反向推论不成立，因为邻接、行优先写回、容量、比例取整和暂态峰值仍可能导致失败。

冷却合成器可使用槽位布尔变量、库存约束、乐观热流和邻接约束。若该模型是所有真实可行布局的必要条件松弛，则 \texttt{UNSAT} 可以安全排除，\texttt{SAT} 只能提交精确模拟。求解器的 \texttt{UNKNOWN} 或超时不得解释为不可行。

当前程序已实现第一组必要条件证书：\texttt{skeleton\_heat\_per\_tick} 精确计算 $Q(s)$；\texttt{sustainable\_vent\_upper\_bound} 计算骨架层乐观持续外排；\texttt{sustainable\_heat\_flow\_upper\_bound} 在完整冷却赋值上建立乐观最大流。该网络保留燃料热源、散热上限、元件散热邻接、抽堆热和换热硬上限，同时放宽行优先顺序、热比例和容量。只有最大流仍小于产热时才排除。

官方换热器在低百分比堆体交换分支可能使用 \texttt{exchange\_side}/2，这可大于名义 \texttt{exchange\_hull}，所以堆体边必须取乐观上界
\[
v_{\mathrm{hull}}^{\max}=\max(\texttt{exchange\_hull},\lfloor\texttt{exchange\_side}/2\rfloor,1).
\]
直接使用名义堆体交换值会低估真实能力，不能用于安全排除。

\begin{theorem}[乐观最大流排除的可靠性]
若完整布局存在安全周期轨迹，则其松弛热流网络的最大流不小于 $Q(s)$。
\end{theorem}

\begin{proof}
在一个完整状态周期上，每个有限储热节点的净热量变化为零。对实际产热、搬运和外排量取时间平均，得到从燃料热源到外界热汇、总流量为 $Q(s)$ 的守恒流。松弛网络保留所有实际通路的硬上限，并只增加有利方向、忽略比例与顺序限制，因此该平均流仍然可行。故若最大流小于 $Q(s)$，安全周期不可能存在。
\end{proof}

\section{前沿动态规划}

令网格为 $G=(V,E)$。对组件类型 $t$ 定义
\[
b(t)=
\begin{cases}
5r_ta_t,&t\text{ 为燃料},\\
0,&\text{其他},
\end{cases}
\]
并对相邻类型 $s,t$ 定义
\[
w(s,t)=5\left[
\mathbf{1}_F(s)r_s\mathbf{1}_A(t)+
\mathbf{1}_F(t)r_t\mathbf{1}_A(s)
\right],
\]
其中 $\mathbf{1}_F$ 表示燃料，$\mathbf{1}_A$ 表示燃料或反射板。则
\[
P(s)=\sum_{v\in V}b(t_v)+\sum_{(u,v)\in E}w(t_u,t_v).
\]

按列逐格扫描高度 $h=6$ 的网格。令 $D(i,F,\mathbf u)$ 为处理前 $i$ 格、前沿标签为 $F$、库存使用量为 $\mathbf u$ 时的最大局部功率。放置新类型 $t$ 时只增加
\[
b(t)+w(t,t_{\mathrm{up}})+w(t,t_{\mathrm{left}}).
\]

\begin{theorem}[前沿动态规划精确性]
每个状态只保留最大前驱，最终值等于库存约束下的 $P_{\max}$，回溯前驱可构造取得该值的骨架。
\end{theorem}

\begin{proof}
对已处理格数归纳。所有从过去连接未来的网格边，其已处理端标签均包含在前沿 $F$ 中；库存剩余量由 $\mathbf u$ 唯一确定。具有相同 $(F,\mathbf u)$ 的较小局部解对任何未来扩展都不会优于较大局部解，因此可安全丢弃。递推枚举每个状态的全部合法新标签，故终态取得全局最大值。
\end{proof}

当前 \texttt{SkeletonPowerTable} 使用精确后缀形式
\[
H(i,F,\mathbf r)=
\max_{t\in A(i,\mathbf r)}
\left[b(t)+w(t,t_{\rm up})+w(t,t_{\rm left})
+H(i+1,F',\mathbf r-\mathbf e_t)\right],
\]
其中 $\mathbf r$ 是剩余库存。终态只有在至少使用一个燃料组件时取零，否则取负无穷。固定发电槽只允许指定标签；功率层空位哨兵 \texttt{POWER\_EMPTY} 在骨架中强制为空，但不占用完整布局槽位，冷却层仍可使用该格。故该值是子问题的精确最大剩余功率，而不是允许库存重复使用的乐观估计。

在总棒数模式下，单联、双联和四联共享一个剩余棒数资源，不分别保存三个无独立上限的使用计数；具有独立上限的反射板仍各占一个资源维度。该约化不改变合法后缀集合。

表以规则版本、结构版本、Python marshal 版本、列数、发电库存、总棒数约束和固定骨架槽为键，压缩持久化到 \texttt{.data/skeleton\_power\_tables.sqlite3}。冷却库存和热学证书不进入功率表键，也不跨镜像共享。

\begin{theorem}[最大堆的骨架功率顺序]
令队列前缀节点的当前功率为 $p$，精确界为 $U=p+H(i,F,\mathbf r)$。反复取最大 $U$ 节点并按下一格全部合法标签分裂，则弹出的完整骨架按功率非增排列。
\end{theorem}

\begin{proof}
子节点前缀集合互斥，且并集恰好等于父节点的全部完整后代。每个子节点的 $U$ 是其后代最大功率。完整节点满足 $U=P$；当它成为堆顶时，任一未决完整骨架都位于某个界不超过 $P$ 的节点中，所以不可能有更高功率骨架尚未输出。归纳即得非增顺序。第一项是最大骨架；继续弹出得到第二个布局，跳过并列值则得到第二个不同功率层。
\end{proof}

Top-$K$ 在 $K>1$ 时只关闭 $U$ 严格低于第 $K$ 名门槛的前缀；Top-1 已有见证后允许关闭 $U$ 等于门槛的前缀。\texttt{count\_power\_subtree} 对互斥堆前缀递推计数最终骨架，并按剩余空位乘 \texttt{count\_cooling\_completions}，所以仍保持完整有标签布局进度恒等式。

当前 \texttt{\_exhaustive\_shards} 按距中心由近到远选择固定位置，只允许发电组件和功率层空位 \texttt{POWER\_EMPTY} 参与前缀；它持续展开全部库存合法功率前缀，直到分片数达到工作进程数的四倍或位置耗尽。

\begin{theorem}[自适应功率前缀分片的互斥与完整性]
设固定位置集合为 $R$。每个完整库存有效布局恰好属于一个自适应功率前缀分片。
\end{theorem}

\begin{proof}
令 $\pi(L)$ 保留燃料与反射板，并把其他标签映射为 \texttt{POWER\_EMPTY}。任意完整布局 $L$ 在 $R$ 上有唯一功率投影 $\pi(L)|_R$。分片器生成全部库存合法功率投影，得到覆盖性；两个不同投影至少在一个固定位置不同，得到互斥性。固定发电组件由第一层扣减；\texttt{POWER\_EMPTY} 只禁止该格放置发电组件，不从第二层冷却空位中扣除，因此每个分片仍完整生成其骨架的全部冷却完成。
\end{proof}

同一骨架及其全部冷却完成只属于一个分片，避免因冷却前缀复制功率表和骨架优先队列。每个分片内部使用精确功率优先队列；进程之间共享排行榜下界，但不承诺跨分片严格同步。该差异不影响完整性。

若骨架标签数为 $k$，库存向量可能值数为 $B\le\prod_t(c_t+1)$，槽位数为 $n=hw$，则
\[
T_{DP}=O(nBk^{h+1}),
\qquad
M_{DP}=O(Bk^h).
\]
直接骨架枚举为 $O(k^n)$。高度固定为 6 时，动态规划消除了骨架上界计算对列数的指数项；这不意味着完整热学可实现性问题也变成多项式时间。

持久化不改变第一次建表复杂度；命中时读取的是实际到达的前沿状态整数表，而不是完整骨架列表。库存维数和上限很大时 $B$ 仍可能很大，故不能预生成从 1 到 216 棒的全部位置布局。

\section{带证书的最优搜索}

算法维护具有精确 DP 功率上界 $U$ 的未决骨架前缀最大堆，以及已经精确验证的可行下界 $P^*$。骨架前缀按下一槽合法标签分裂。每次处理最大 $U$ 的前缀，完整骨架依次应用有限反射板、热守恒、直接冷却构造和必要条件热流；无法证明的分支回退到完整冷却生成和精确模拟。当最大未决上界满足适用于请求 Top-$K$ 的排行榜闭合条件时停止。

\begin{theorem}[构造式求解器全局最优性]
若 DP 上界精确，所有排除证书可靠，所有可行见证通过精确模拟，并且未被证明排除的分支最终被验证或保持未决，则在 $U_{\max}\le P^*$ 时返回的布局是边界 $T$ 下的全局最优 Mark I。
\end{theorem}

\begin{proof}
任意尚未验证的可行布局属于某个未决子问题，其功率不超过该子问题上界，故不超过 $U_{\max}\le P^*$。已关闭分支由可靠证书证明不可行或不可能超过下界。返回布局已经精确验证且取得 $P^*$，所以不存在更高功率的可行布局。
\end{proof}

\begin{remark}[完备性边界]
固定点、热守恒和静态热流约束都不是完整的 Mark I 判定。全局完备性来自精确回退。若关闭回退，结果只能标记为当前最佳证书布局，不得设置 \texttt{proven\_global}。
\end{remark}

\section{新增算法的效率边界}

令 $N$ 为完整库存有效布局数，$M$ 为新算法精确模拟的不同布局数，$B_c$ 为由目标上界、数学证书或可靠 \texttt{UNSAT} 结论关闭而无需模拟的布局数。在实现不重复模拟且保留完整回退时，概念上
\[
N=M+B_c,
\qquad M\le N.
\]
若 $B_c>0$，则精确模拟次数严格减少。设一次模拟平均成本为 $C_t$，动态规划、约束求解和调度额外成本为 $C_o$，相对完整模拟的严格墙钟加速条件为
\[
(N-M)C_t>C_o.
\]
因此可以证明 DP 的骨架上界复杂度改进和精确模拟次数不增加；不能无条件证明每个输入的总墙钟时间严格缩短。最坏情况下证书均无效，算法仍可能回退到组合搜索。

只有满足以下条件才可报告全局证明：DP 状态包含完整前沿与库存；上界不低估；\texttt{SAT} 候选精确复核；\texttt{UNKNOWN} 不关闭分支；no-good 只排除已有证明的赋值；镜像不共享热学证书；停止时不存在上界高于 $P^*$ 的未决分支。

第一阶段实现基准使用 3 列、单铀一个、铱反射板一个、普通散热片一个、单工作进程和 40000 反应堆周期边界。完整空间为 5526；加入证书前执行 5526 次精确模拟、约 70.9 秒，加入热守恒与乐观热流证书后执行 802 次模拟并证明排除 4724 个布局、约 1.0 秒，二者最优值均为 5 EU/t。该数据证明此输入上的模拟次数减少约 85.5\%，不构成对所有输入的统一加速保证。

冷却见证构造器的 Top-1 进程内隔离基准使用 3 列、单铀一个和普通散热片上限四个，共 57852 个完整布局。关闭构造器约 18.1 ms，开启后约 0.78 ms；两者均最终验证 1 个、计数跳过 57851 个且最优值均为 5 EU/t。约 23 倍差异来自避免递归展开同骨架冷却排列，不代表复杂换热库存或 Top-$K$ 的统一倍数。

功率层分片结构测试使用 3 列、单铀上限 4 和 31 个工作进程，在第 8 个中心优先位置得到 163 个互斥功率前缀，且不含冷却标签。只有一根单铀时，功率空间至多产生 19 个分片；程序不再为凑满核心而把同一骨架的冷却完成拆到不同进程。

持久功率表隔离基准使用 9 列、最多四个单铀、Top-1 和单进程。旧骨架路径共有 342540 个完整骨架，在无候选基准约 63.1 秒。配置二十个反应堆散热片时，当前精确表包含 6272 个状态，首次建表并证明 60 EU/t 见证约 0.033 秒，再次磁盘命中约 0.008 秒，只精确验证一个布局。无持续散热组件时，根级最低产热/最大散热比较不建 DP 表，约 0.001 秒证明全部 342540 个骨架不可实现。后一倍数不能全部归因于 DP；大量并列高功率骨架均难以判定时仍可能产生较多队列节点。

\section{程序映射}

\begin{longtable}{p{0.34\textwidth}p{0.56\textwidth}}
\toprule
数学对象或定理 & 程序实现 \\
\midrule
约化状态 $S$ & \texttt{ReactorSimulator.state\_signature} \\
固定点快进 & \texttt{ReactorSimulator.\_fast\_forward\_fixed\_state} \\
发电骨架 $P(L)$ & \texttt{power\_skeleton} \\
骨架功率 & \texttt{skeleton\_eu\_per\_tick} \\
二层双射枚举 & \texttt{\_run\_mark\_i\_two\_level\_shard} \\
冷却子树组合数 & \texttt{count\_cooling\_completions} \\
有限反射板引理 & \texttt{has\_degrading\_power\_component} \\
多进程门槛 & \texttt{shared\_power\_floor} \\
进程级结果证书 & \texttt{\_fixed\_point\_certificate} \\
启发式重复布局缓存 & \texttt{OptimizationJob.\_heuristic\_cache} \\
骨架精确产热 & \texttt{skeleton\_heat\_per\_tick} \\
持续散热上界 & \texttt{sustainable\_vent\_upper\_bound} \\
乐观持续热流 & \texttt{sustainable\_heat\_flow\_upper\_bound} \\
精确后缀功率表 & \texttt{SkeletonPowerTable.best\_suffix} \\
持久骨架表 & \texttt{.data/skeleton\_power\_tables.sqlite3} \\
功率降序骨架 & \texttt{SkeletonPowerTable.ranked\_skeletons} 与分片内最大堆 \\
功率前缀完整计数 & \texttt{count\_power\_subtree}（二层生成器内部） \\
自适应功率层分片 & \texttt{\_exhaustive\_shards(..., power\_only=True, target\_shards=...)} \\
拟实现：冷却约束合成 & 必要条件松弛；\texttt{SAT} 后调用精确模拟 \\
最优停止证书 & 精确队列界闭合且未决分支为空 \\
\bottomrule
\end{longtable}

\end{document}
